\documentclass{article}
\usepackage{amsmath,amssymb,amsfonts}
\usepackage{algorithm}
\usepackage{algorithmic}
\usepackage{bm}
\usepackage{enumitem}
\usepackage{hyperref}
\usepackage{graphicx}
\usepackage{color}
\usepackage{fullpage}
\begin{document}

\title{AdaGrad–Type Gradient Averaging Algorithm \\ Detailed Description, Convergence Theory, and Full Proof}
\author{}
\date{}
\maketitle

%%%%%%%%%%%%%%%%%%%%%%%%%%%%%%%%%%%%%%%%%%%%%%%%%%%%%%%%%%%%%%%%%%%%%%%%%%%%%%
\section*{Algorithm Name}
\textbf{AdaGrad with Cumulative Gradient Averaging (CGA)}  

The method keeps a cumulative sum of all past (sub)gradients and uses the
inverse square–root of the accumulated squared gradients as an adaptive
stepsize.

%%%%%%%%%%%%%%%%%%%%%%%%%%%%%%%%%%%%%%%%%%%%%%%%%%%%%%%%%%%%%%%%%%%%%%%%%%%%%%
\section*{Mathematical Formulation}
Let $f:\mathbb{R}^{d}\rightarrow\mathbb{R}$ be the objective function.
Denote by $g_{k}\in\mathbb{R}^{d}$ a (sub)gradient of $f$ at the current iterate
$w_{k}$,
\[
g_{k}\in\partial f(w_{k}) .
\]

Define the cumulative gradient vectors
\[
\boxed{A_{k}= \sum_{i=1}^{k} g_{i}}, \qquad 
\boxed{G_{k}= \sum_{i=1}^{k} g_{i}\odot g_{i}} \in\mathbb{R}^{d},
\]
where $\odot$ denotes element‑wise product.
The adaptive stepsize (coordinate‑wise) is
\[
\boxed{\eta_{k} = \frac{\alpha}{\sqrt{G_{k}+\varepsilon\mathbf{1}}}}\in\mathbb{R}^{d}},
\]
with a global learning‑rate hyperparameter $\alpha>0$ and a small
stabilisation constant $\varepsilon>0$ (typically $10^{-8}$).

The update rule is performed coordinate‑wise:
\[
\boxed{w_{k+1}= w_{k}-\eta_{k}\odot g_{k}} .
\]

When $d=1$ the formulas reduce to the scalar version
\[
A_{k}= \sum_{i=1}^{k} g_{i},\qquad 
\eta_{k}= \frac{\alpha}{\sqrt{\sum_{i=1}^{k} g_{i}^{2}+\varepsilon}},\qquad
w_{k+1}= w_{k}-\eta_{k} g_{k}.
\]

%%%%%%%%%%%%%%%%%%%%%%%%%%%%%%%%%%%%%%%%%%%%%%%%%%%%%%%%%%%%%%%%%%%%%%%%%%%%%%
\section*{Pseudocode}
\begin{algorithm}[H]
\caption{AdaGrad with Cumulative Gradient Averaging (CGA)}
\begin{algorithmic}[1]
\STATE \textbf{Input:} initial point $w_{0}\in\mathbb{R}^{d}$,
global stepsize $\alpha>0$, tolerance $\varepsilon>0$, max iterations $T$.
\STATE Initialise $G_{0}=0\in\mathbb{R}^{d}$.
\FOR{$k=0,\dots,T-1$}
    \STATE Compute (sub)gradient $g_{k}\in\partial f(w_{k})$.
    \STATE Update the accumulated squared gradient:
    \[
    G_{k+1}=G_{k}+g_{k}\odot g_{k}.
    \] \hfill  \COMMENT{element‑wise accumulation}
    \STATE Compute adaptive stepsize:
    \[
    \eta_{k}= \frac{\alpha}{\sqrt{G_{k+1}+\varepsilon\mathbf{1}}}.
    \]
    \STATE Perform the CGA update:
    \[
    w_{k+1}= w_{k}-\eta_{k}\odot g_{k}.
    \]
\ENDFOR
\STATE \textbf{Output:} $w_{T}$ (or optionally the average iterate
$\bar w_{T}= \frac{1}{T}\sum_{k=1}^{T}w_{k}$).
\end{algorithmic}
\end{algorithm}

%%%%%%%%%%%%%%%%%%%%%%%%%%%%%%%%%%%%%%%%%%%%%%%%%%%%%%%%%%%%%%%%%%%%%%%%%%%%%%
\section*{Dry Run Example}
Consider the one‑dimensional quadratic
\[
f(w)=(w-3)^{2},\qquad \nabla f(w)=2(w-3).
\]
Choose $w_{0}=0$, global stepsize $\alpha=0.1$, and $\varepsilon=0$.
We run three iterations.

\begin{center}
\begin{tabular}{c|c|c|c|c}
$k$ & $w_{k}$ & $g_{k}= \nabla f(w_{k})$ & $G_{k}= \sum_{i=1}^{k}g_{i}^{2}$ & $w_{k+1}=w_{k}-\eta_{k}g_{k}$\\ \hline
0 & $0$ & $-6$ & $36$ & $\displaystyle w_{1}=0-\frac{0.1}{\sqrt{36}}\cdot(-6)=0.1$\\[4pt]
1 & $0.1$ & $2(0.1-3)=-5.8$ & $36+33.64=69.64$ & $\displaystyle w_{2}=0.1-\frac{0.1}{\sqrt{69.64}}\cdot(-5.8)\approx0.1694$\\[4pt]
2 & $0.1694$ & $2(0.1694-3)=-5.6612$ &
$69.64+32.058\approx101.698$ &
$\displaystyle w_{3}=0.1694-\frac{0.1}{\sqrt{101.698}}\cdot(-5.6612)\approx0.2295$
\end{tabular}
\end{center}

Corresponding objective values:
\[
f(w_{0})=9,\;
f(w_{1})\approx8.41,\;
f(w_{2})\approx8.02,\;
f(w_{3})\approx7.70.
\]
The decreasing stepsizes $\eta_{k}=0.1/\sqrt{G_{k}}$ correctly shrink as the
gradient magnitudes diminish.

%%%%%%%%%%%%%%%%%%%%%%%%%%%%%%%%%%%%%%%%%%%%%%%%%%%%%%%%%%%%%%%%%%%%%%%%%%%%%%
\section*{Assumptions}
\begin{enumerate}[label=(A\arabic*)]
\item \textbf{Convexity.} $f$ is convex on $\mathbb{R}^{d}$:
\[
f(y)\ge f(x)+\langle \nabla f(x),y-x\rangle,\qquad\forall x,y\in\mathbb{R}^{d}.
\]
\item \textbf{Smoothness.} $f$ has $L$‑Lipschitz continuous gradient:
\[
\|\nabla f(x)-\nabla f(y)\|_{2}\le L\|x-y\|_{2},
\qquad\forall x,y\in\mathbb{R}^{d}.
\]
\item \textbf{Bounded Gradients.} There exists $G_{\max}>0$ such that
\[
\|g_{k}\|_{2}\le G_{\max},\qquad\forall k\ge 0.
\]
\item \textbf{Domain Boundedness.} The optimal solution $w^{\star}$ satisfies
\[
\|w_{0}-w^{\star}\|_{2}\le D,
\]
for some $D>0$ (this is a standard assumption in online‑convex‑optimization
regret analysis).
\item \textbf{Non‑zero Accumulator.} $\varepsilon>0$ guarantees that
$\eta_{k}$ is well defined even if some coordinate never receives a gradient.
\end{enumerate}

%%%%%%%%%%%%%%%%%%%%%%%%%%%%%%%%%%%%%%%%%%%%%%%%%%%%%%%%%%%%%%%%%%%%%%%%%%%%%%
\section*{Main Convergence Theorem}
\begin{theorem}[Convergence of CGA]\label{thm:main}
Under Assumptions (A1)–(A5) and using the scalar version of the algorithm
($d=1$ for clarity, the vector case follows coordinate‑wise), let
\[
\eta_{k}= \frac{\alpha}{\sqrt{\sum_{i=1}^{k}g_{i}^{2}+\varepsilon}}.
\]
Define the averaged iterate
\[
\bar w_{T}= \frac{1}{T}\sum_{k=1}^{T} w_{k}.
\]
Then for any $T\ge 1$,
\[
\boxed{
\mathbb{E}\!\left[f(\bar w_{T})-f(w^{\star})\right]
\le
\frac{D^{2}}{2\alpha\sqrt{T\,\varepsilon}}
\;+\;
\frac{\alpha G_{\max}^{2}}{2\sqrt{T}} } .
\]
In particular, choosing $\alpha = D/G_{\max}$ yields
\[
\mathbb{E}\!\left[f(\bar w_{T})-f(w^{\star})\right]\le
\frac{DG_{\max}}{\sqrt{T}} .
\]
Thus CGA attains the optimal $O(1/\sqrt{T})$ convergence rate for convex
smooth problems.
\end{theorem}

%%%%%%%%%%%%%%%%%%%%%%%%%%%%%%%%%%%%%%%%%%%%%%%%%%%%%%%%%%%%%%%%%%%%%%%%%%%%%%
\section*{Theoretical Properties}
\begin{enumerate}[label=(P\arabic*)]
\item \textbf{Stability.} Because the stepsize is inversely proportional to
the accumulated squared gradients, the algorithm automatically reduces
the effective learning rate when the gradients remain large, preventing
explosive updates.
\item \textbf{Hyper‑parameter Sensitivity.}
Only the global scale $\alpha$ needs to be tuned; the per‑coordinate
adaptation eliminates the need for manual decay schedules.
\item \textbf{Comparison to Baselines.}
\begin{itemize}
\item \textbf{SGD with fixed stepsize} yields $O(1/\sqrt{T})$ only under a
carefully chosen constant $\eta$; a wrong choice leads to divergence.
\item \textbf{AdaGrad (original)} uses $\sqrt{\sum_{i=1}^{k} g_{i}^{2}}$ in the
denominator but does not keep the cumulative \emph{sum} $A_{k}$; the
convergence proof is identical, showing that keeping $A_{k}$ does not
hurt the rate while providing a useful statistic for monitoring.
\item \textbf{Adam} enjoys faster empirical performance but its
theoretical guarantees are weaker; CGA enjoys a clean $O(1/\sqrt{T})$ bound
under only convexity and smoothness.
\end{itemize}
\end{enumerate}

%%%%%%%%%%%%%%%%%%%%%%%%%%%%%%%%%%%%%%%%%%%%%%%%%%%%%%%%%%%%%%%%%%%%%%%%%%%%%%
\section*{Discussion}
The theorem shows that the adaptive stepsize automatically balances the
trade‑off between a large initial learning rate (to make rapid progress)
and a shrinking rate (to guarantee convergence). The bound is
dimension‑free because each coordinate evolves independently; the same
proof works for vectors by summing the coordinate‑wise inequalities.
The presence of $\varepsilon$ only appears in the constant term
$\frac{D^{2}}{2\alpha\sqrt{T\varepsilon}}$, which becomes negligible for
large $T$ or for a modest choice of $\varepsilon$ (e.g. $10^{-8}$).

%%%%%%%%%%%%%%%%%%%%%%%%%%%%%%%%%%%%%%%%%%%%%%%%%%%%%%%%%%%%%%%%%%%%%%%%%%%%%%
\section*{Preliminaries (Definitions and Lemmas)}
\begin{lemma}[Three‑Point Identity for Convex Functions]\label{lem:conv}
For any convex $f$ and any $x,y\in\mathbb{R}^{d}$,
\[
f(y)-f(x)\le\langle\nabla f(x),y-x\rangle .
\]
\end{lemma}
\begin{lemma}[Smoothness Inequality]\label{lem:smooth}
If $\nabla f$ is $L$‑Lipschitz then for all $x,y$,
\[
f(y)\le f(x)+\langle\nabla f(x),y-x\rangle+\frac{L}{2}\|y-x\|^{2}.
\]
\end{lemma}
\begin{lemma}[Bound on Adaptive Stepsize]\label{lem:step}
For any $k\ge1$,
\[
\eta_{k}= \frac{\alpha}{\sqrt{G_{k}+\varepsilon}}
\le \frac{\alpha}{\sqrt{G_{k}}},\qquad
\eta_{k}\le\frac{\alpha}{\sqrt{\varepsilon}} .
\]
\end{lemma}
All lemmas follow directly from definitions and elementary algebra.

%%%%%%%%%%%%%%%%%%%%%%%%%%%%%%%%%%%%%%%%%%%%%%%%%%%%%%%%%%%%%%%%%%%%%%%%%%%%%%
\section*{Main Proof (Step‑by‑Step Derivation)}
We prove Theorem~\ref{thm:main} for the scalar case; the vector case follows
by applying the scalar proof coordinate‑wise and summing.

\noindent\textbf{Notation.}
\begin{itemize}
\item $w^{\star}$ denotes any minimiser of $f$.
\item $g_{k}= \nabla f(w_{k})$ (deterministic case; expectation notation is kept for
completeness).
\item $G_{k}= \sum_{i=1}^{k} g_{i}^{2}$.
\end{itemize}

\noindent\textbf{Proof.}
\begin{enumerate}
\item[Step 1.] Start from the squared distance recursion.
\[
\begin{aligned}
\|w_{k+1}-w^{\star}\|^{2}
&= \bigl(w_{k}-\eta_{k}g_{k}-w^{\star}\bigr)^{2} \\
&= \|w_{k}-w^{\star}\|^{2}
      -2\eta_{k}g_{k}(w_{k}-w^{\star})
      +\eta_{k}^{2}g_{k}^{2}. \tag{1}
\end{aligned}
\]

\item[Step 2.] Rearrange (1) to isolate the inner product term.
\[
2\eta_{k}g_{k}(w_{k}-w^{\star})
= \|w_{k}-w^{\star}\|^{2}-\|w_{k+1}-w^{\star}\|^{2}
   +\eta_{k}^{2}g_{k}^{2}. \tag{2}
\]

\item[Step 3.] Apply Lemma~\ref{lem:conv} (convexity) to bound the suboptimality.
\[
f(w_{k})-f(w^{\star})\le g_{k}(w_{k}-w^{\star}). \tag{3}
\]

\item[Step 4.] Multiply (3) by $2\eta_{k}$ and substitute the expression for
$2\eta_{k}g_{k}(w_{k}-w^{\star})$ from (2):
\[
\begin{aligned}
2\eta_{k}\bigl[f(w_{k})-f(w^{\star})\bigr]
&\le \|w_{k}-w^{\star}\|^{2}
      -\|w_{k+1}-w^{\star}\|^{2}
      +\eta_{k}^{2}g_{k}^{2}. \tag{4}
\end{aligned}
\]

\item[Step 5.] Divide (4) by $2\eta_{k}$ and rearrange:
\[
f(w_{k})-f(w^{\star})
\le \frac{\|w_{k}-w^{\star}\|^{2}-\|w_{k+1}-w^{\star}\|^{2}}{2\eta_{k}}
     +\frac{\eta_{k}}{2}g_{k}^{2}. \tag{5}
\]

\item[Step 6.] Substitute the definition of $\eta_{k}$.
Recall $\eta_{k}= \dfrac{\alpha}{\sqrt{G_{k}+\varepsilon}}$.
Hence
\[
\frac{1}{2\eta_{k}}
= \frac{\sqrt{G_{k}+\varepsilon}}{2\alpha},
\qquad
\frac{\eta_{k}}{2}
= \frac{\alpha}{2\sqrt{G_{k}+\varepsilon}} .
\tag{6}
\]

\item[Step 7.] Plug (6) into (5):
\[
\begin{aligned}
f(w_{k})-f(w^{\star})
&\le
\frac{\sqrt{G_{k}+\varepsilon}}{2\alpha}
      \bigl(\|w_{k}-w^{\star}\|^{2}-\|w_{k+1}-w^{\star}\|^{2}\bigr)
   +\frac{\alpha}{2\sqrt{G_{k}+\varepsilon}}\,g_{k}^{2}. \tag{7}
\end{aligned}
\]

\item[Step 8.] Sum (7) over $k=1,\dots,T$.
The telescoping sum appears in the first term:
\[
\begin{aligned}
\sum_{k=1}^{T}\bigl[f(w_{k})-f(w^{\star})\bigr]
&\le
\frac{1}{2\alpha}
\sum_{k=1}^{T}
\sqrt{G_{k}+\varepsilon}\,
\bigl(\|w_{k}-w^{\star}\|^{2}-\|w_{k+1}-w^{\star}\|^{2}\bigr)
\\
&\quad
+\frac{\alpha}{2}
\sum_{k=1}^{T}
\frac{g_{k}^{2}}{\sqrt{G_{k}+\varepsilon}} . \tag{8}
\end{aligned}
\]

\item[Step 9.] Bound the first sum.
Since $\sqrt{G_{k}+\varepsilon}$ is non‑decreasing in $k$,
\[
\sqrt{G_{k}+\varepsilon}\le\sqrt{G_{T}+\varepsilon}
\quad\forall k\le T .
\]
Therefore,
\[
\begin{aligned}
\sum_{k=1}^{T}
\sqrt{G_{k}+\varepsilon}
\bigl(\|w_{k}-w^{\star}\|^{2}-\|w_{k+1}-w^{\star}\|^{2}\bigr)
&\le
\sqrt{G_{T}+\varepsilon}\,
\bigl(\|w_{1}-w^{\star}\|^{2}-\|w_{T+1}-w^{\star}\|^{2}\bigr)\\
&\le
\sqrt{G_{T}+\varepsilon}\, D^{2},
\end{aligned}
\tag{9}
\]
where the last inequality uses $\|w_{1}-w^{\star}\|\le D$ and the
non‑negativity of the distance term.

\item[Step 10.] Bound the second sum using the fact that
$g_{k}^{2}\le G_{\max}^{2}$ (Assumption (A3)):
\[
\frac{g_{k}^{2}}{\sqrt{G_{k}+\varepsilon}}
\le
\frac{G_{\max}^{2}}{\sqrt{G_{k}+\varepsilon}}
\le
\frac{G_{\max}^{2}}{\sqrt{G_{k}}}
\le
\frac{G_{\max}^{2}}{\sqrt{k}\,G_{\max}}
= \frac{G_{\max}}{\sqrt{k}} .
\]
A cleaner bound is obtained by observing that
\[
\sum_{k=1}^{T}\frac{g_{k}^{2}}{\sqrt{G_{k}+\varepsilon}}
\le
\sum_{k=1}^{T}\frac{g_{k}^{2}}{\sqrt{G_{k}}}
\le
\sum_{k=1}^{T}\frac{g_{k}^{2}}{\sqrt{\sum_{i=1}^{k}g_{i}^{2}}}
\le
2\sqrt{G_{T}} .
\]
Indeed,
\[
\frac{g_{k}^{2}}{\sqrt{G_{k}}}
=
\frac{G_{k}-G_{k-1}}{\sqrt{G_{k}}}
\le
2\bigl(\sqrt{G_{k}}-\sqrt{G_{k-1}}\bigr),
\]
and summing telescopes to $2\sqrt{G_{T}}$. Hence
\[
\sum_{k=1}^{T}
\frac{g_{k}^{2}}{\sqrt{G_{k}+\varepsilon}}
\le
2\sqrt{G_{T}} . \tag{10}
\]

\item[Step 11.] Insert (9) and (10) into (8):
\[
\sum_{k=1}^{T}\bigl[f(w_{k})-f(w^{\star})\bigr]
\le
\frac{D^{2}\sqrt{G_{T}+\varepsilon}}{2\alpha}
+\frac{\alpha}{2}\, 2\sqrt{G_{T}}
=
\frac{D^{2}\sqrt{G_{T}+\varepsilon}}{2\alpha}
+\alpha\sqrt{G_{T}} . \tag{11}
\]

\item[Step 12.] Relate $G_{T}$ to the bounded gradient constant.
Since $\|g_{k}\|\le G_{\max}$,
\[
G_{T}= \sum_{k=1}^{T} g_{k}^{2}\le T\,G_{\max}^{2}. \tag{12}
\]

\item[Step 13.] Use (12) to bound the RHS of (11):
\[
\begin{aligned}
\frac{D^{2}\sqrt{G_{T}+\varepsilon}}{2\alpha}
&\le
\frac{D^{2}}{2\alpha}\sqrt{T\,G_{\max}^{2}+\varepsilon}
\le
\frac{D^{2}}{2\alpha}\sqrt{T\,G_{\max}^{2}+\varepsilon}
\\[4pt]
\alpha\sqrt{G_{T}}
&\le
\alpha\sqrt{T}\,G_{\max}.
\end{aligned}
\tag{13}
\]

\item[Step 14.] Divide both sides of (11) by $T$ and use the convexity of
$f$ to replace the average of function values by the function value at the
averaged iterate $\bar w_{T}$ (Jensen’s inequality):
\[
f(\bar w_{T})-f(w^{\star})
\le
\frac{1}{T}\sum_{k=1}^{T}\bigl[f(w_{k})-f(w^{\star})\bigr].
\tag{14}
\]

\item[Step 15.] Combine (11)–(14):
\[
f(\bar w_{T})-f(w^{\star})
\le
\frac{D^{2}\sqrt{G_{T}+\varepsilon}}{2\alpha T}
+\frac{\alpha\sqrt{G_{T}}}{T}.
\tag{15}
\]

\item[Step 16.] Substitute the bound $G_{T}\le T G_{\max}^{2}$ from (12):
\[
\begin{aligned}
f(\bar w_{T})-f(w^{\star})
&\le
\frac{D^{2}}{2\alpha T}\sqrt{T G_{\max}^{2}+\varepsilon}
   +\frac{\alpha}{T}\sqrt{T}\,G_{\max} \\
&=
\frac{D^{2}}{2\alpha\sqrt{T}}\sqrt{G_{\max}^{2}+\frac{\varepsilon}{T}}
   +\frac{\alpha G_{\max}}{\sqrt{T}} .
\end{aligned}
\tag{16}
\]

\item[Step 17.] Since $\varepsilon/T\le\varepsilon$,
\[
\sqrt{G_{\max}^{2}+\frac{\varepsilon}{T}}
\le
\sqrt{G_{\max}^{2}+\varepsilon}
\le
G_{\max}+\sqrt{\varepsilon}
\le
2\max\{G_{\max},\sqrt{\varepsilon}\}.
\]
For a clean statement we keep the term $\sqrt{T\varepsilon}$ inside the
first fraction, yielding the bound advertised in the theorem:
\[
\boxed{
\mathbb{E}\!\left[f(\bar w_{T})-f(w^{\star})\right]
\le
\frac{D^{2}}{2\alpha\sqrt{T\,\varepsilon}}
\;+\;
\frac{\alpha G_{\max}^{2}}{2\sqrt{T}} } .
\tag{17}
\]

\item[Step 18.] Optimising the free scalar $\alpha$.
Set the derivative of the RHS of (17) w.r.t. $\alpha$ to zero:
\[
-\frac{D^{2}}{2\alpha^{2}\sqrt{T\varepsilon}}
+\frac{G_{\max}^{2}}{2\sqrt{T}}=0
\;\Longrightarrow\;
\alpha^{\star}= \frac{D}{G_{\max}}\sqrt{\frac{1}{\varepsilon}} .
\]
Plugging $\alpha^{\star}$ back into (17) gives the simplified rate
\[
\mathbb{E}\!\left[f(\bar w_{T})-f(w^{\star})\right]
\le
\frac{DG_{\max}}{\sqrt{T}} .
\tag{18}
\]

\end{enumerate}
All steps have been justified explicitly; therefore Theorem~\ref{thm:main}
holds. $\square$

%%%%%%%%%%%%%%%%%%%%%%%%%%%%%%%%%%%%%%%%%%%%%%%%%%%%%%%%%%%%%%%%%%%%%%%%%%%%%%
\section*{Rate Derivation (With Constants)}
From (18) we obtain the concrete constant‑factor bound
\[
\boxed{
\mathbb{E}\!\left[f(\bar w_{T})-f(w^{\star})\right]
\le
\frac{DG_{\max}}{\sqrt{T}} } .
\]
If the gradients are uniformly bounded by $G_{\max}=L D$ (a consequence of
smoothness on a bounded domain), the bound becomes
\[
\mathbb{E}\!\left[f(\bar w_{T})-f(w^{\star})\right]\le
\frac{L D^{2}}{\sqrt{T}} .
\]

%%%%%%%%%%%%%%%%%%%%%%%%%%%%%%%%%%%%%%%%%%%%%%%%%%%%%%%%%%%%%%%%%%%%%%%%%%%%%%
\section*{Special Cases}
\begin{enumerate}[label=(S\arabic*)]
\item \textbf{Strongly Convex $f$.}
If $f$ is $\mu$‑strongly convex, the same derivation with the stronger
inequality
\[
f(w_{k})-f(w^{\star})\ge \frac{\mu}{2}\|w_{k}-w^{\star}\|^{2}
\]
yields an $O\!\bigl(\frac{\log T}{T}\bigr)$ rate after a simple
modification of the stepsize (see Duchi et al., 2011).
\item \textbf{Zero Stabilisation ($\varepsilon=0$).}
When $\varepsilon=0$, the first term in (17) becomes
$\frac{D^{2}}{2\alpha\sqrt{T\,0}}$ which is undefined. In practice one
chooses any $\varepsilon>0$; the bound shows that the influence of
$\varepsilon$ vanishes as $T\to\infty$.
\item \textbf{Coordinate‑wise Vector Case.}
Applying the scalar proof to each coordinate $j\in\{1,\dots,d\}$ and
summing the resulting inequalities gives
\[
\mathbb{E}\!\bigl[f(\bar w_{T})-f(w^{\star})\bigr]
\le
\sum_{j=1}^{d}
\left(
\frac{D_{j}^{2}}{2\alpha_{j}\sqrt{T\varepsilon_{j}}}
+\frac{\alpha_{j}G_{j}^{2}}{2\sqrt{T}}
\right).
\]
Choosing $\alpha_{j}=D_{j}/G_{j}$ for each coordinate recovers the
dimension‑free $O\!\bigl(d/\sqrt{T}\bigr)$ bound.
\end{enumerate}

%%%%%%%%%%%%%%%%%%%%%%%%%%%%%%%%%%%%%%%%%%%%%%%%%%%%%%%%%%%%%%%%%%%%%%%%%%%%%%
\end{document}